% Generated by tools/build.js from content/resume.yaml -- edit the YAML, not this file.
\documentclass[10pt,letterpaper]{article}
\usepackage[T1]{fontenc}
\usepackage[utf8]{inputenc}
\usepackage{lmodern}
\usepackage{textcomp}
\usepackage[margin=0.6in]{geometry}
\usepackage{enumitem}
\usepackage[hidelinks]{hyperref}
\hypersetup{pdftitle={Jesús M. Rueda-Becerril, Ph.D.}, pdfauthor={Jesús M. Rueda-Becerril}}
% map every glyph (including ligatures) to Unicode so text extraction is exact
\input{glyphtounicode}
\pdfgentounicode=1

\pagestyle{empty}
\setlength{\parindent}{0pt}
\setlength{\parskip}{0pt}
\setlength{\emergencystretch}{3em}
\setlist[itemize]{leftmargin=1.3em, topsep=1pt, itemsep=0.5pt, parsep=0pt}
\newcommand{\cvsection}[1]{%
  \vspace{5pt}{\large\bfseries #1}\par\nopagebreak\vspace{2pt}\hrule
  \nopagebreak\vspace{3pt}\nopagebreak}
\newcommand{\cvsubsection}[1]{%
  \vspace{3pt}{\bfseries #1}\par\nopagebreak\vspace{2pt}\nopagebreak}
\newcommand{\entrygap}{\vspace{3pt}}
\newcommand{\leftright}[2]{{#1\unskip\nobreak\hfil\penalty50\hskip1em\hbox{}%
  \nobreak\hfil\mbox{#2}\parfillskip=0pt\par}}

\begin{document}

\begin{center}
  {\LARGE\bfseries Jesús M. Rueda-Becerril, Ph.D.}\\[3pt]
  Seattle, WA \textbar{} \href{mailto:jm.ruebe@gmail.com}{jm.ruebe@gmail.com} \textbar{} +1 (765) 430-2330 \textbar{} \href{https://linkedin.com/in/jeruebe}{linkedin.com/in/jeruebe} \textbar{} \href{https://altjerue.github.io}{altjerue.github.io}
\end{center}

\cvsection{Professional Summary}
Computational Scientist with a PhD in Physics and over ten years of experience building and running high-performance numerical algorithms and physics-based simulations. Expertise in scientific machine learning and geospatial data analysis. Strong mathematical foundations with a track record of architecting open-source simulation libraries, leading funded research projects, and translating complex physical models into production-quality scientific software. Experienced working across cross-functional research and engineering teams.\par

\cvsection{Skills}
\textbf{Languages:} C/C++, Python, Fortran 95, Rust, R, Java, SQL, Julia, Shell\par
\textbf{ML \& Data:} Scientific machine learning, PyTorch, TensorFlow, scikit-learn, NumPy, SciPy, Pandas, Astropy, ML-assisted parameter estimation, gradient-based optimization, SVM, random forest, neural networks\par
\textbf{HPC \& Parallel:} OpenMP, OpenACC (GPU offloading), MPI, SLURM, PBS, HDF5\par
\textbf{Tools:} Git, Docker, Azure, Jenkins, Kafka, Jira, QGIS, ParaView, VisIt\par

\cvsection{Experience}
\leftright{\textbf{Independent Contractor}}{Jan 2026 – Present}\nopagebreak
\leftright{\textit{Mercor}}{Remote}\nopagebreak
Physics subject-matter expert evaluating frontier language models for a leading AI lab.\par\nopagebreak
\begin{itemize}
  \item{} Reviewed 50+ expert-level problems in physics, hydrodynamics, astrophysics, and cosmology; wrote original prompts that stumped frontier models.
\end{itemize}
\entrygap

\leftright{\textbf{Spatial Data Scientist}}{May 2025 – Jan 2026}\nopagebreak
\leftright{\textit{TealWaters}}{Seattle, WA}\nopagebreak
Built geospatial ML pipelines and optimized terrain-analysis code for wetland mapping.\par\nopagebreak
\begin{itemize}
  \item{} Optimized the terrain-modeling code that computes elevation derivatives from 1 m, 3 m, and 10 m DEMs over 20 tiles, improving its computational performance.
  \item{} Mapped wetland probability across the 2,160 km² Skykomish watershed with the Wetland Intrinsic Potential (WIP) random-forest model, reaching >75\% accuracy.
\end{itemize}
\entrygap

\leftright{\textbf{Independent Researcher \& Open-Source Developer}}{Jan 2024 – May 2025}\nopagebreak
\leftright{\textit{Self-directed}}{Seattle, WA}\nopagebreak
Continued simulation research and open-source development between industry roles.\par\nopagebreak
\begin{itemize}
  \item{} Co-developed Tleco, a Rust/Python toolkit that solves the Fokker–Planck equation coupled to radiative transfer for relativistic outflows (ApJ 2024).
  \item{} Developing WindsOfChange with the UW Forest Mycobiome Lab: a C++/R library for wind-driven spore dispersal over hilly forest terrain (RcppParallel, OpenMP).
\end{itemize}
\entrygap

\leftright{\textbf{Software Engineer}}{Apr 2022 – Jan 2024}\nopagebreak
\leftright{\textit{Paychex}}{Remote}\nopagebreak
Built and tested data-integration services between enterprise systems and Oracle E-Business Suite.\par\nopagebreak
\begin{itemize}
  \item{} Developed and tested production data-pipeline software (Java, Python, SQL) with CI/CD.
\end{itemize}
\entrygap

\leftright{\textbf{Postdoctoral Research Associate}}{Feb 2021 – Apr 2022}\nopagebreak
\leftright{\textit{Rochester Institute of Technology}}{Rochester, NY}\nopagebreak
HPC simulations of supermassive black hole binaries in an NSF-funded project.\par\nopagebreak
\begin{itemize}
  \item{} Led a 7-person team in an NSF-sponsored upgrade of PatchworkMHD, a C code for MHD simulations of supermassive black hole binaries, adding black hole spin with under 5\% runtime overhead.
  \item{} Ran and benchmarked large production simulations on Frontera (TACC), about 10,000 node-hours, and identified performance bottlenecks in the parallel code.
  \item{} Calibrated Paramo parameters against Fermi-LAT observations using gradient-based optimization, comparing simulation output with data and updating the model automatically; supervised a graduate student on the same problem.
  \item{} Collaborated on a multi-institutional HPC project on binary neutron star mergers, resulting in 2 publications.
\end{itemize}
\entrygap

\leftright{\textbf{Postdoctoral Research Fellow}}{Oct 2018 – Nov 2020}\nopagebreak
\leftright{\textit{Purdue University}}{West Lafayette, IN}\nopagebreak
Radiative-transfer modeling of relativistic jets (blazars and gamma-ray burst afterglows).\par\nopagebreak
\begin{itemize}
  \item{} Designed and wrote Paramo, an open-source Fortran 95 library that solves the Fokker–Planck equation for relativistic particle distributions coupled to radiative transfer. OpenMP gives about 20× speedup on 5 cores; GPU offload tested with OpenACC.
  \item{} Implemented implicit PDE solvers, numerical integration, and arbitrary-shape particle distributions without slowing the simulation.
  \item{} Wrote Python tools to analyze Paramo output and run parameter optimization; other groups have adopted Paramo, including Combi \& Siegel (2023).
  \item{} Co-I on two NASA Fermi Guest Investigator proposals, one funded (\$68K) to model blazar radiation; mentored 3 graduate students and published 2 peer-reviewed papers.
\end{itemize}
\entrygap

\leftright{\textbf{Postdoctoral Research Fellow}}{Jan 2018 – Sep 2018}\nopagebreak
\leftright{\textit{Universidad Michoacana de San Nicolás de Hidalgo}}{Morelia, Mexico}\nopagebreak
\begin{itemize}
  \item{} Developed a Python pipeline to process black hole simulation outputs from the GRTrans numerical code, generating SVM training data for radio image prediction.
  \item{} Built open-source data analysis and visualization tools in Python for radiative transfer phenomena in relativistic astrophysics scenarios.
\end{itemize}
\entrygap

\leftright{\textbf{Graduate Research Assistant}}{Oct 2011 – Jul 2017}\nopagebreak
\leftright{\textit{Universitat de València}}{Valencia, Spain}\nopagebreak
\begin{itemize}
  \item{} Studied how magnetization of shocked plasma shapes the radiation spectra of relativistic internal shocks, using the C-SPEV simulation code.
\end{itemize}
\entrygap


\cvsection{Projects}
\leftright{\textbf{Blazar Population ML}}{Jan 2026 – Present}\nopagebreak
\leftright{\textit{Scientific ML (Python, scikit-learn, PyTorch)}}{\href{https://github.com/altjerue/ilwikak-ml}{github.com/altjerue/ilwikak-ml}}\nopagebreak
\begin{itemize}
  \item{} Applying classical ML and Bayesian inference to the Fermi-LAT 4LAC-DR3 catalog (3,407 sources) to recover physical jet parameters from multi-wavelength data.
  \item{} Building a neural network that maps physical jet parameters to observable spectral energy distributions accurately enough to compare with Tleco results.
  \item{} Forecasting blazar gamma-ray flares with a time-series foundation model fine-tuned on Tleco-generated synthetic light curves.
\end{itemize}
\entrygap


\cvsection{Selected Publications (9 peer-reviewed articles, 3 first-author)}
\begin{itemize}
  \item{} Davis, Z., \textbf{Rueda-Becerril, J. M.}, \& Giannios, D. \textit{Tleco: A Toolkit for Modeling Radiative Signatures from Relativistic Outflows}, \href{https://doi.org/10.3847/1538-4357/ad8bc2}{ApJ 976, 182 (2024)}, \href{https://arxiv.org/abs/2405.17581}{arXiv:2405.17581}.
  \item{} \textbf{Rueda-Becerril, J. M.}, Harrison, A. O., \& Giannios, D. \textit{Blazar jets launched with similar energy per baryon, independently of their power}, \href{https://doi.org/10.1093/mnras/staa3925}{MNRAS 501, 4092–4102 (2021)}, \href{https://arxiv.org/abs/2009.02273}{arXiv:2009.02273}.
  \item{} \textbf{Rueda-Becerril, J. M.}, Mimica, P., \& Aloy, M. A. \textit{On the influence of a hybrid thermal–non-thermal distribution in the internal shocks model for blazars}, \href{https://doi.org/10.1093/mnras/stx476}{MNRAS 468, 1169–1182 (2017)}, \href{https://arxiv.org/abs/1612.06383}{arXiv:1612.06383}.
\end{itemize}

\cvsection{Education}
\leftright{\textbf{Universitat de València}}{Valencia, Spain}\nopagebreak
\leftright{\textit{Ph.D. in Physics, Excellent Cum Laude}}{2017}
\entrygap

\leftright{\textbf{Universidad Michoacana de San Nicolás de Hidalgo}}{Morelia, Mexico}\nopagebreak
\textit{M.Sc. in Physics}\par
\entrygap

\leftright{\textbf{Universidad Autónoma del Estado de México}}{Toluca, Mexico}\nopagebreak
\textit{B.Sc. in Physics}\par
\entrygap


\end{document}
